% BEGIN PROFESSOR FEEDBACK PF06: Foundations - Temporal motifs and fingerprint
\section{Temporal motifs}
\label{sec:motifs}

\paragraph{Static network motifs.}
The concept of network motifs was introduced by \citet{milo2002network} to describe small, recurring subgraph patterns that appear significantly more often in a real network than in a randomised reference. These patterns are considered the building blocks of network structure, so that different classes of networks (biological, social, technological) tend to be characterised by different sets of over-represented motifs \citep{milo2002network}. In static graphs, a motif is formally defined as a connected induced subgraph on $k$ nodes whose frequency in the observed network exceeds what would be expected by chance under a suitable null model \citep{milo2002network, newman2010networks}.


\paragraph{$\delta$-temporal motifs.}
\citet{paranjape2017motifs} extend the static motif concept to temporal networks by introducing $\delta$-temporal motifs.
A $k$-node, $l$-edge $\delta$-temporal motif is an ordered sequence of
$l$ timestamped edges involving $k$ nodes. The graph formed by these edges must
be connected, and the time between the first and last edge must be at most
$\delta$. Two occurrences represent the same motif when their nodes can be
relabelled without changing the direction or temporal order of the edges
\citep{paranjape2017motifs}.

\paragraph{Candidate patterns and statistical motifs.}
The temporal counting framework defines patterns by their structure, edge order,
and time span. Counting an occurrence does not itself establish that the pattern
is more common than expected by chance. In this thesis, an untested pattern is
therefore a \emph{candidate motif}, or a temporal subgraph. The terms
\emph{motif role} and \emph{motif fingerprint} refer to the features derived from
these candidates. Statistical overrepresentation would require a separate
comparison with randomised networks under a specified null model
\citep{milo2002network}. The supervised tests used here assess prediction and
category separability, not network overrepresentation.

\paragraph{The catalogue of $\delta$-temporal motifs.}

For motifs with three edges on two or three nodes,
\citet{paranjape2017motifs} enumerate a complete catalogue of 36 patterns.
It contains four \emph{2-node motifs}, in which all three edges occur between
the same pair of nodes, 24 \emph{star motifs}, in which one centre node
interacts with two distinct neighbours, and eight \emph{triangle motifs}, in
which every pair of the three nodes is connected. Only two of the eight
triangle motifs form a directed cycle. The individual patterns also differ in
the direction and temporal ordering of their edges. The full catalogue is
illustrated in Figure~\ref{fig:motifs_paranjape}.

\begin{figure}[h]
    \centering
    \includegraphics[width=\textwidth]{motifs_paranjape}
    \caption{All 36 two- and three-node, three-edge $\delta$-temporal motifs, 
    indexed as $M_{i,j}$ by row $i$ and column $j$. The first edge in each 
    motif originates from the green node and edge labels denote temporal ordering. 
    Figure from \citet{paranjape2017motifs}.}
    \label{fig:motifs_paranjape}
\end{figure}

\subsection{Counting algorithm}

Motif counting algorithms form the computational foundation of any motif-based analysis. This section first outlines the established methods for static graphs and then explains why a fundamentally different approach is required for temporal data.

\subsubsection{Static motif counting}

Static graph motif counting is a well-studied problem. Given a graph $G$, the goal is to count how many times a given subgraph pattern, such as a triangle or a directed cycle, appears in $G$. Algorithms for this task range from exact neighborhood-intersection methods to matrix-based approaches and sampling-based approximations, covering a broad spectrum of accuracy and efficiency trade-offs \citep{Yin2025}. These methods are computationally mature and underpin many downstream tasks such as clustering coefficient computation and community detection.

However, static methods operate on a time-aggregate representation: all edges that ever existed between two nodes collapse into a single connection, regardless of when they occurred \citep{kovanen2011temporal}. This means that the temporal ordering is lost entirely. \citet{paranjape2017motifs} show that even when a triangle is topologically present in a network, only a subset of its instances are temporally valid. The others either occur in a different order or fall outside a given time window $\delta$. A static counter treats all of them as identical. For blockchain transaction networks, this distinction is critical because coordinated malicious behaviour happens through transaction sequences that have a specific temporal order, something that static aggregation cannot recover.

\subsubsection{Temporal motif counting}
Temporal motif counting extends the static counting problem by requiring that all edges of a motif occur in a specific order within a bounded time window $\delta$. 

The general counting algorithm proceeds in three steps: (1) identify all
static instances of the subgraph $H$ in the static graph $G$,
(2) for each instance, gather all corresponding temporal edges into a
time-sorted sequence $S$, and (3) process $S$ within a sliding window of
size $\delta$, tracking how often each ordered edge pattern appears
within that window \citep{paranjape2017motifs}. Algorithm~\ref{alg:paranjape}
formalizes this procedure.

\begin{figure}[t]
\begin{minipage}{\linewidth}

\begin{algorithm}[H]
\caption{$\delta$-Temporal Motif Counting adapted from \citet{paranjape2017motifs}}
\label{alg:paranjape}
\begin{algorithmic}[1]
\Require Time-sorted sequence $S = (e_1, \ldots, e_L)$,
         window $\delta$, motif length $l$
\Ensure Count of each $l$-edge $\delta$-temporal motif in $S$
\State $\textit{left} \gets 1$
\State $\textit{patternCounts} \gets \emptyset$
\For{$\textit{right} = 1, \ldots, L$}
    \While{$t_{\textit{right}} - t_{\textit{left}} > \delta$}
        \State \Call{RemoveFromWindow}{$e_{\textit{left}}$}
        \State $\textit{left} \gets \textit{left} + 1$
    \EndWhile
    \State \Call{AddToWindow}{$e_{\textit{right}}$}
\EndFor
\State \Return $\textit{patternCounts}$
\end{algorithmic}

\noindent\rule{\linewidth}{0.4pt}
\noindent\textbf{Procedure} \textsc{RemoveFromWindow}$(e)$
\hfill\textit{// expire $e$ and withdraw prefix contributions}
\begin{algorithmic}[1]
\State $\textit{patternCounts}[e] \gets \textit{patternCounts}[e] - 1$
\ForAll{suffix $\sigma$ of length $< l-1$ in $\textit{patternCounts}$}
    \State $\textit{patternCounts}[(e,\,\sigma)]
           \mathrel{-}= \textit{patternCounts}[\sigma]$
\EndFor
\end{algorithmic}

\noindent\rule{\linewidth}{0.4pt}
\noindent\textbf{Procedure} \textsc{AddToWindow}$(e)$
\hfill\textit{// admit $e$ and extend all existing prefixes}
\begin{algorithmic}[1]
\ForAll{prefix $\pi$ of length $< l$ in $\textit{patternCounts}$,
        longest first}
    \State $\textit{patternCounts}[(\pi,\,e)]
           \mathrel{+}= \textit{patternCounts}[\pi]$
\EndFor
\State $\textit{patternCounts}[e] \gets \textit{patternCounts}[e] + 1$
\end{algorithmic}

\end{algorithm}
\end{minipage}
\end{figure}
 
For 2-node motifs this runs in 
$O(m)$ time, which is optimal. For 3-node, 3-edge star and triangle motifs, specialized algorithms reduce the worst-case complexity significantly, with the fast triangle algorithm achieving up to a 56.5$\times$ speedup over the general method on the Bitcoin transaction network \citep{paranjape2017motifs}. This scalability is relevant to the present work, as blockchain transaction graphs are large and the motifs we are interested in correspond directly to 3-node star and triangle patterns within a short time window.

\subsection{Motif fingerprint}
The identification of recurring subgraph patterns in a network provides insight
into the local structural properties of individual nodes. A global count records
how often a pattern occurs anywhere in the graph. It does not say which position
an individual node occupies. The \emph{motif-role fingerprint} of
\citet{McDonnell2014} retains this distinction by collecting counts for the local
roles of a focal node, such as an initiator, intermediary, or receiver.

\citet{McDonnell2014} further show that functional motif-role fingerprints, obtained from counts of partial subgraphs, can be transformed into structural fingerprints by multiplying them with an invertible transformation matrix, establishing a mathematically well-defined correspondence between the two representations. The authors identify anomaly detection and the identification of structurally similar nodes as primary use cases, which makes it well applicable for this thesis.

The concept is extended here to temporal graphs following the
$\delta$-temporal counting framework. A global candidate pattern first defines
the ordered transfers. Choosing a focal address then determines its local role
in that occurrence. Counts of these roles form the address fingerprint.
Table~\ref{tab:motif_terminology} distinguishes these levels.

\begin{table}[htbp]
\centering
\small
\caption{The terms used to describe temporal patterns and address features.
The counted patterns are candidates, without a claim of statistical
overrepresentation.}
\label{tab:motif_terminology}
\begin{tabular}{@{}p{2.9cm} p{3.6cm} p{6.1cm}@{}}
\toprule
\textbf{Term} & \textbf{Example} & \textbf{Meaning} \\
\midrule
Global candidate motif &
$A \to B$, then $A \to C$, then $A \to B$ &
The complete ordered pattern, considered without selecting a focal address.
A statistical motif additionally requires evidence of overrepresentation
against a network null model. \\[5pt]
Focal node &
Address $A$ &
The address whose participation is being described. \\[5pt]
Local motif role &
$A$ as the outgoing star centre &
The position of the focal address, including the directions and ordering of
the transfers relative to it. \\[5pt]
Motif class &
Star, all outgoing &
A structural grouping of related patterns or roles. A class is not one
feature column. \\[5pt]
Address fingerprint &
$\mathbf{f}_v^\delta$ &
The complete vector of local role counts for address $v$, with one component
for each retained role. \\
\bottomrule
\end{tabular}
\end{table}

For the full local catalogue used in this thesis, the raw fingerprint is
defined as

\begin{equation}
    \mathbf{f}_v^{\delta}
      = \bigl(c_{v,1}^{\delta},\,c_{v,2}^{\delta},\,\ldots,\,c_{v,40}^{\delta}\bigr)
    \label{eq:motif_fingerprint}
\end{equation}

where $c_{v,j}^{\delta}$ counts the number of candidate occurrences in which
address $v$ participates in local role $j$ within a span of at most $\delta$.
The 40 components are local motif-role counts, not 40 global patterns or 40
motif classes. Raphtory reports 24 star-centre roles, eight two-node roles, and
eight triangle roles. Section~\ref{sec:motifcountingandretainedroles} explains
their relationship to the 36 global patterns and the 32-role subsets used when
triangle coverage is incomplete. A feature matrix has one row per address and
one column per retained local role.

This serves two purposes. First, it reduces the high-dimensional transaction graph to a compact, interpretable, and fixed-size representation per address. Second, it encodes behavioral structure rather than raw transaction statistics because an address involved in coordinated activity is expected to exhibit a distinct fingerprint compared to ordinary users, as coordinated behavior by definition produces recurring, time-constrained interaction patterns among a group of addresses. The fingerprint therefore creates a behavioral signature of each address showing small recurring temporal subgraph patterns.
% END PROFESSOR FEEDBACK PF06
